% P16-C3 and a finite identification component of P16-C4.
\chapter{Finite Structural Causal Models}
\label{sec:p16-finite-causal-models}

A structural causal model specifies how its variables are assigned values
and how those assignments change under a specified intervention. Pearl's
nonparametric structural equations and equation-replacement semantics supply
the source presentation used here
\cite[secs.~3.1--3.2.3, pp.~103--112]{Pearl2009CausalOverview}.
This chapter gives a finite formulation with complete proofs, including the
probability assumptions required for factorization. All numerical examples
are stipulated mathematical models.

The probability language is that of
\cref{sec:p15-finite-conditioning,sec:p15-finite-independence}.
The background implications of logical abduction in
\cref{sec:p16-finite-abduction} do not themselves specify the assignments
or interventions defined below.

\section{Graphs, assignments and disturbances}
\label{sec:p16-causal-model-data}

\begin{hbdefinition}{def:p16-causal-dag}{A finite directed acyclic graph}
Let \(N=\{1,\ldots,n\}\), with \(n\geq1\). A \emph{directed graph}
\(G=(N,E)\) consists of the vertices \(N\) and a set
\(E\subseteq N\times N\) of directed edges. Write \(j\to i\)
for \((j,i)\in E\), and put
\[
 \operatorname{pa}_G(i)\defeq\{j\in N:j\to i\}.
\]
These vertices are the \emph{parents} of \(i\). A directed path of
positive length is a sequence \(i_0,\ldots,i_m\), with \(m\geq1\),
such that \(i_{r-1}\to i_r\) for every \(r\). A directed cycle
is such a path with \(i_0=i_m\). The graph is \emph{acyclic} when
there is no directed cycle, including a loop of length one. A finite
directed acyclic graph is abbreviated \emph{DAG}.

For \(I\subseteq N\), let \(\operatorname{de}_G(I)\) be the vertices
reachable from a member of \(I\) by a directed path of positive length.
A \emph{topological order} is an enumeration \(i_1,\ldots,i_n\) of
all vertices in which every parent precedes its child.
\index{causal graph!directed acyclic graph}\index{topological order}
\end{hbdefinition}

\begin{hbtheorem}{thm:p16-causal-topological-order}{proved=annex:p16.causal-topological-order}{Topological evaluation is available}
Every finite DAG has a topological order. Repeatedly removing a vertex
with no incoming edge in the remaining graph produces one after exactly
\(n\) removals.
\end{hbtheorem}

\begin{hbdefinition}{def:p16-finite-scm}{A finite structural causal model}
A model \(M\) consists of the following data.
\begin{enumerate}
\item A DAG \(G=(N,E)\), whose vertices index the \emph{endogenous}
      variables to be assigned by the equations, and a finite nonempty
      value set \(S_i\) for each vertex \(i\). For \(J\subseteq N\), put
      \(S_J=\prod_{j\in J}S_j\). Products are indexed tuples;
      \(S_\varnothing\) contains the single empty tuple. If
      \(v\in S_N\), write \(v_J=(v_j)_{j\in J}\).
\item A finite nonempty outcome space \(\Omega\), with a mass
      function \(q:\Omega\to[0,1]\) satisfying
      \(\sum_{\omega\in\Omega}q(\omega)=1\). Its event probability
      is \(\mathbb P(A)=\sum_{\omega\in A}q(\omega)\).
\item For each \(i\), a finite nonempty disturbance value set
      \(\mathcal U_i\), a function \(U_i:\Omega\to\mathcal U_i\),
      and a total structural function
      \[
       f_i:S_{\operatorname{pa}_G(i)}\times\mathcal U_i\longrightarrow S_i.
      \]
\end{enumerate}
The \(U_i\) are \emph{exogenous disturbances}: this model gives their
joint probability law but no equations assigning them from the endogenous
variables. They may share information or be dependent. Their independence
will be an additional assumption when explicitly stated.

At an outcome \(\omega\), a \emph{solution} is an assignment
\(v\in S_N\) satisfying all the structural equations
\[
 v_i=f_i(v_{\operatorname{pa}_G(i)},U_i(\omega)),\qquad i\in N.
\]
The indexed left sides declare which variable each function assigns.
Solving an equation algebraically for another variable does not replace
that declaration. The parent sets state which endogenous inputs a function
is allowed to use. A listed parent may be ignored by the function; an edge
does not assert a nonzero effect or a detectable association.
\index{structural causal model!finite}\index{exogenous disturbance}
\end{hbdefinition}

\begin{hbtheorem}{thm:p16-causal-solution}{proved=annex:p16.causal-solution}{Existence, uniqueness and the observational law}
For every \(\omega\in\Omega\), the model has exactly one solution.
Evaluating the structural functions in any topological order gives that
solution. Denote its coordinates by \(V_i(\omega)\), and its tuple by
\(V(\omega)\). The \(V_i\) are finite random variables. Their joint
\emph{observational law} is the mass function
\[
 p_M(v)\defeq\mathbb P(V=v)
       =\sum_{\omega:V(\omega)=v}q(\omega),\qquad v\in S_N.
\]
For \(J\subseteq N\), its marginal law is
\[
 p_{M,J}(s)=\mathbb P(V_J=s)
           =\sum_{v:v_J=s}p_M(v),\qquad s\in S_J.
\]
All these laws are defined even when some outcomes or assignments have
probability zero.
\end{hbtheorem}

\section{Interventions replace specified equations}
\label{sec:p16-finite-interventions}

\begin{hbdefinition}{def:p16-causal-intervention}{A fixed intervention and its distribution}
Choose \(I\subseteq N\) and \(x\in S_I\). The intervention
\(\operatorname{do}(V_I=x)\) replaces the equation for every \(i\in I\)
by \(v_i=x_i\). It preserves \(\Omega,q\), all disturbances and
every equation indexed outside \(I\). Write \(M_{I=x}\) for this
model. Its graph deletes every edge whose terminal vertex is in \(I\).
For \(i\in I\), the new structural function is the constant \(x_i\),
with no endogenous parent; its old disturbance can be retained as an
unused input.

This operation stipulates that the unreplaced mechanisms and the
exogenous probability law remain invariant under the modeled manipulation.
That stipulation is part of the causal interpretation of the model. If an
actual manipulation also changes another mechanism or the population's
exogenous law, the displayed replacement alone does not model that change.

Write \(V^{I=x}(\omega)\) for its solution and
\(p_M^{I=x}(v)=\mathbb P(V^{I=x}=v)\) for its law. The traditional
intervention notation means
\[
 \mathbb P_M(V_J=s\mid\operatorname{do}(V_I=x))
 \defeq\mathbb P(V_J^{I=x}=s).
\]
Here the expression following the vertical bar names an operation on the
model. It is not a conditioning event. In particular this definition
requires no assumption that \(\mathbb P(V_I=x)>0\).

For an event \(D\subseteq S_J\), sum the displayed law over \(s\in D\).
For a real-valued function \(h:S_J\to\mathbb R\), the corresponding
intervention mean is the finite expectation
\(\mathbb E[h(V_J^{I=x})]\). A difference of such means at two fixed
interventions is a model's \emph{average intervention effect} for that
specified response function and contrast.
\index{intervention!equation replacement}\index{do operator}
\end{hbdefinition}

\begin{hbtheorem}{thm:p16-causal-intervention-laws}{proved=annex:p16.causal-intervention-laws}{Laws of fixed equation replacement}
Every intervention above has a unique solution at every outcome, and
\[
 p_M^{I=x}(v)=\sum_{\omega:V^{I=x}(\omega)=v}q(\omega)
\]
is a mass function. It vanishes unless \(v_I=x\). In addition:
\begin{enumerate}
\item If \(V_I(\omega)=x\), then
      \(V^{I=x}(\omega)=V(\omega)\).
\item If \(j\notin I\cup\operatorname{de}_G(I)\), then
      \(V_j^{I=x}(\omega)=V_j(\omega)\) for every \(\omega\).
\item Intervening on the empty set leaves the model unchanged.
      For disjoint \(I,J\), performing the two replacements
      \(I=x\) and \(J=y\) in either order gives the same model
      as the single replacement on \(I\cup J\) with these values.
\end{enumerate}
The first two conclusions are pointwise, including outcomes of zero mass.
\end{hbtheorem}

Conditioning on \(V_I=x\), when this event has positive probability,
reweights the original outcomes to
\(q(\omega)\mathbf1_{\{V_I=x\}}(\omega)/\mathbb P(V_I=x)\)
and retains the original solution functions. Equation replacement retains
the original weights and evaluates the modified functions. The confounded
example below calculates the resulting difference explicitly.

\section{Independent disturbances and factorization}
\label{sec:p16-causal-factorization}

\begin{hbdefinition}{def:p16-causal-kernels}{Independent disturbances and structural kernels}
Put \(r_i(u)=\mathbb P(U_i=u)\). The disturbances are
\emph{mutually independent} if
\[
 \mathbb P(U_1=u_1,\ldots,U_n=u_n)=\prod_{i=1}^n r_i(u_i)
 \quad\text{for every }(u_i)_{i\in N}\in\prod_{i\in N}\mathcal U_i.
\]
This is independence of the entire family, not just of every pair.
Following the source, a model with an acyclic graph and mutually
independent disturbances is called \emph{Markovian}
\cite[sec.~3.2.3, p.~110]{Pearl2009CausalOverview}.

For any model, define its structural kernels by
\[
 k_i(b\mid a)\defeq
 \sum_{\substack{u\in\mathcal U_i\\f_i(a,u)=b}}r_i(u),
 \qquad a\in S_{\operatorname{pa}_G(i)},\quad b\in S_i.
\]
For each fixed \(a\), this is a mass function on \(S_i\): the sets
of disturbance values giving each output partition \(\mathcal U_i\).
The kernel is defined at every parent assignment, including assignments
which never occur observationally. Its vertical bar separates output from
input; it does not define a conditional probability on a null event.
\index{causal model!Markovian}\index{structural kernel}
\end{hbdefinition}

\begin{hbtheorem}{thm:p16-causal-factorization}{proved=annex:p16.causal-factorization}{The finite causal factorization}
In a model with mutually independent disturbances,
\[
 p_M(v)=\prod_{i\in N}k_i(v_i\mid v_{\operatorname{pa}_G(i)})
 \qquad(v\in S_N).
\]
If \(\mathbb P(V_{\operatorname{pa}_G(i)}=a)>0\), then
\[
 k_i(b\mid a)=
 \mathbb P(V_i=b\mid V_{\operatorname{pa}_G(i)}=a).
\]
Thus the structural kernels supply the observational conditional factors
on every parent event on which conditioning is defined. The product formula
itself requires no positivity assumption.
\end{hbtheorem}

\begin{hbtheorem}{thm:p16-causal-truncation}{proved=annex:p16.causal-truncation}{Truncated factorization for simultaneous interventions}
In the same model, for \(I\subseteq N\), \(x\in S_I\), and
\(v\in S_N\),
\[
 p_M^{I=x}(v)=
 \begin{cases}
 \displaystyle\prod_{i\notin I}k_i(v_i\mid v_{\operatorname{pa}_G(i)}),
       &v_I=x,\\[1ex]
 0,    &v_I\ne x.
 \end{cases}
\]
In the first case the values in intervened parent positions are therefore
the specified components of \(x\). Marginal intervention probabilities
are obtained by summing this mass function over the remaining coordinates.
For \(I=N\) the empty product is one, giving point mass at \(x\).
\end{hbtheorem}

This is a fully defined finite version of the source's causal Markov
condition and truncated factorization
\cite[Theorem~1 and Corollary~1, pp.~110--111]{Pearl2009CausalOverview}.
Using observational conditional probabilities in place of the kernels
also requires that the relevant conditioning events have positive mass.
An intervention can assign a parent configuration that was absent from
the observational distribution.

\section{Identification from a specified observational law}
\label{sec:p16-finite-causal-identification}

\begin{hbdefinition}{def:p16-causal-identification}{Identification relative to a model class}
Let \(\mathcal C\) be a class of models with common observed coordinates
\(O\) and value sets \(S_O\), and let \(Q(M)\) be a specified
quantity associated with each model in the class. The quantity is
\emph{identified from the observational law on \(O\)} within
\(\mathcal C\) if, for every \(M_1,M_2\in\mathcal C\),
\[
 p_{M_1,O}=p_{M_2,O}\quad\Longrightarrow\quad Q(M_1)=Q(M_2).
\]
Equality of laws means equality at every assignment in \(S_O\).
A pair with equal observed laws and unequal values of \(Q\) proves
\emph{nonidentification} in that class. This is the two-model criterion
of the source's Definition~2 \cite[p.~109]{Pearl2009CausalOverview}.
Identification treats the observational law as given exactly; recovering
it or the target from finite sampled data is a separate estimation problem.
\index{causal identification}\index{nonidentification}
\end{hbdefinition}

\begin{hbtheorem}{thm:p16-causal-positive-identification}{proved=annex:p16.causal-positive-identification}{Identification with full observation and positive support}
Fix the DAG and finite value sets. Consider the class of models with
mutually independent disturbances in which all endogenous variables are
observed and \(p_M(v)>0\) for every \(v\in S_N\).
For every fixed intervention \(I=x\), its full distribution is identified
in this class. Explicitly, for \(v_I=x\), it is
\[
 p_M^{I=x}(v)=\prod_{i\notin I}
 \frac{\displaystyle\sum_{w:w_i=v_i,\,
                      w_{\operatorname{pa}_G(i)}=v_{\operatorname{pa}_G(i)}}p_M(w)}
      {\displaystyle\sum_{w:w_{\operatorname{pa}_G(i)}=v_{\operatorname{pa}_G(i)}}p_M(w)}.
\]
It is zero for \(v_I\ne x\). Every denominator is positive.
Consequently all its marginals and finite response means are identified.
Full positive support is a sufficient hypothesis here; no claim that it is
necessary for every particular target is made.
\end{hbtheorem}

\begin{hbexample}{def:p16-causal-chain}{The source's three-variable chain in the finite formulation}
Let the endogenous variables be \(Z,X,Y\), with graph
\(Z\to X\to Y\) and mutually independent disturbances. Write their
kernels as \(k_Z(z),k_X(x\mid z),k_Y(y\mid x)\).
Then the preceding theorems give
\[
\begin{aligned}
 p_M(z,x,y)&=k_Z(z)k_X(x\mid z)k_Y(y\mid x),\\
 \mathbb P_M(Z=z,Y=y\mid\operatorname{do}(X=x))
     &=k_Z(z)k_Y(y\mid x).
\end{aligned}
\]
Summing the second expression over \(z\) gives \(k_Y(y\mid x)\),
because \(\sum_z k_Z(z)=1\). If \(\mathbb P(X=x)>0\), this equals
\(\mathbb P(Y=y\mid X=x)\) by
\cref{thm:p16-causal-factorization}. This specializes the source's
equations~(5)--(7) and (16)--(17). The stated independence and graph,
together with positive mass for the conditional expression, justify the
coincidence in this example.
\end{hbexample}

\section{A confounded association and two indistinguishable observations}
\label{sec:p16-causal-confounding-example}

\begin{hbdefinition}{def:p16-confounded-binary-model}{A binary common-cause model}
For bits \(a,b\in\{0,1\}\), let \(a\oplus b\) equal zero when
\(a=b\) and one otherwise. Take
\(\Omega=\{0,1\}^3\), with coordinates \((z,a,b)\), and mass
\[
 q(z,a,b)=\tfrac12\alpha(a)\alpha(b),\qquad
 \alpha(0)=\tfrac34,\quad\alpha(1)=\tfrac14.
\]
The coordinate disturbances \(U_Z(z,a,b)=z\),
\(U_X(z,a,b)=a\), and \(U_Y(z,a,b)=b\) are mutually independent.
Specify endogenous variables with values \(\{0,1\}\) by
\[
 Z=U_Z,\qquad X=Z\oplus U_X,\qquad Y=Z\oplus U_Y.
\]
Their DAG has edges \(Z\to X\) and \(Z\to Y\).
Call this model \(M_A\). The shared parent \(Z\) is a common cause:
its equation feeds both assignments. The following calculation makes the
confounding of the observed \(X,Y\) association precise in this model.
\end{hbdefinition}

\begin{hbtheorem}{thm:p16-causal-confounded-example}{proved=annex:p16.causal-confounded-example}{Association differs from intervention in the common-cause model}
The observational marginal of \((X,Y)\) in \(M_A\) is
\[
\begin{array}{c|cc|c}
 &Y=0&Y=1&\text{row total}\\\hline
 X=0&5/16&3/16&1/2\\
 X=1&3/16&5/16&1/2\\\hline
 \text{column total}&1/2&1/2&1
\end{array}
\]
Hence
\[
 \mathbb P(Y=1\mid X=1)=\tfrac58,\qquad
 \mathbb P(Y=1\mid X=0)=\tfrac38.
\]
For either \(x\in\{0,1\}\), however,
\[
 \mathbb P_{M_A}(Y=1\mid\operatorname{do}(X=x))=\tfrac12.
\]
With response \(Y\), the difference of the two observational conditional
means is \(1/4\), whereas the average intervention effect of setting
\(X=1\) instead of \(X=0\) is zero.

If \(Z\) is also observed, every \((X,Z)\) assignment has positive
mass and this model satisfies the explicit adjustment identity
\[
 \mathbb P_{M_A}(Y=1\mid\operatorname{do}(X=x))
 =\sum_{z=0}^1\mathbb P(Y=1\mid X=x,Z=z)\mathbb P(Z=z).
\]
Its two response factors are \(1/4\) for \(z=0\) and \(3/4\)
for \(z=1\). The identity here is proved directly from this model.
\end{hbtheorem}

\begin{hbtheorem}{thm:p16-causal-nonidentification}{proved=annex:p16.causal-nonidentification}{Equal observational laws with unequal intervention laws}
The following two nonidentification statements hold.
\begin{enumerate}
\item Among finite acyclic models of binary observed variables \(X,Y\)
      with allowed parent edge \(X\to Y\) and possibly dependent
      disturbances, the full observational law does not identify
      \(\mathbb P_M(Y=1\mid\operatorname{do}(X=1))\), even when
      every observed pair has positive probability. There are two such
      models with the table above as their common law, and target values
      \(1/2\) and \(5/8\).
\item In the class with the same binary graph and mutually independent
      disturbances, removing the positive-support hypothesis also permits
      nonidentification. There are two such models with observational
      point mass at \((X,Y)=(0,0)\), but target values zero and one
      under \(\operatorname{do}(X=1)\).
\end{enumerate}
\end{hbtheorem}

\begin{hbwarning}{def:p16-causal-latent-projection}{Dependent disturbances do not inherit the product formula}
The first common-cause example can be expressed using only endogenous
\(X,Y\): on the same \(\Omega\), set
\(\widetilde U_X=(z,a)\), \(\widetilde U_Y=(z,b)\),
\(X=z\oplus a\), and \(Y=z\oplus b\). Allow the parent edge
\(X\to Y\) while the function for \(Y\) ignores that input.
The two disturbances both contain \(z\), and are dependent: their
first coordinates agree with probability one, though each is a fair bit.
Each structural output kernel is the constant mass function \((1/2,1/2)\).
Multiplying those kernels would give \(1/4\) at every pair, contradicting
the proved \(5/16,3/16\) table. The general structural model and its
interventions remain defined; the independent-disturbance hypothesis of
\cref{thm:p16-causal-factorization} is absent.
\end{hbwarning}

\begin{hbmethod}{def:p16-causal-finite-calculation}{Calculating a specified finite intervention}
List the graph, value sets, outcome weights, disturbance functions and
structural functions. Verify that every function is total on its stated
finite domain and that the graph is acyclic. Choose the intervention and
replace its equations. Evaluate once at each \(\omega\) in a
topological order, and add \(q(\omega)\) to the mass of the resulting
assignment. Group these masses for the requested marginal or response
mean. \Cref{thm:p16-causal-solution,thm:p16-causal-intervention-laws}
prove that this finite procedure terminates and returns the specified law.

With mutually independent disturbances, one may instead calculate the
kernels and apply \cref{thm:p16-causal-truncation}. To recover kernels
from observational data, separately establish the conditional identity
and its positive denominators. State the model class when claiming an
identified effect. A table fitting observed frequencies does not provide
the structural equations, their stability under intervention, the
independence of disturbances, or measurement of an omitted variable.
\end{hbmethod}

These structural foundations establish finite acyclic equation-replacement
semantics, factorization under independent disturbances, a sufficient
identification result, and explicit nonidentification examples.
The following sections develop graphical separation and back-door
adjustment in \cref{sec:p16-graphical-separation-adjustment}, then
front-door identification in \cref{sec:p16-front-door-identification}.
Each result states its graph and support assumptions and supplies its
own proof. \Cref{sec:p16-finite-counterfactuals} develops the finite
same-unit counterfactual calculus and the binary potential-outcome bridge.
General counterfactual identification and cyclic-model semantics require
their further constructions.
